\documentclass[10pt]{article}

\usepackage[a4paper,margin=1.75cm]{geometry}
\usepackage[T1]{fontenc}
\usepackage{mathpazo}
\usepackage{microtype}
\usepackage{booktabs}
\usepackage{tabularx}
\usepackage{array}
\usepackage{graphicx}
\usepackage{float}
\usepackage{xcolor}
\usepackage{titlesec}
\usepackage{caption}
\usepackage{enumitem}
\usepackage{fancyhdr}
\usepackage{hyperref}
\usepackage{xurl}

\definecolor{navy}{HTML}{17324D}
\definecolor{blue}{HTML}{2A6F97}
\definecolor{teal}{HTML}{2A9D8F}
\definecolor{purple}{HTML}{7B2CBF}
\definecolor{orange}{HTML}{D97706}
\definecolor{green}{HTML}{198754}
\definecolor{red}{HTML}{B42318}
\definecolor{ink}{HTML}{243B53}
\definecolor{muted}{HTML}{627D98}
\definecolor{paperblue}{HTML}{EEF7FB}
\definecolor{papergreen}{HTML}{EEF8F3}
\definecolor{paperorange}{HTML}{FFF7E8}

\hypersetup{
  colorlinks=true,
  linkcolor=blue,
  urlcolor=purple,
  citecolor=blue,
  pdftitle={Cross-model transfer of evolved tabular FeatureGraphs},
  pdfauthor={GigaEvo}
}
\titleformat{\section}{\Large\bfseries\color{navy}}{\thesection}{0.65em}{}
\titleformat{\subsection}{\large\bfseries\color{blue}}{\thesubsection}{0.55em}{}
\captionsetup{font=small,labelfont={bf,color=navy}}
\setlength{\parindent}{0pt}
\setlength{\parskip}{5pt}
\setlength{\tabcolsep}{5pt}
\renewcommand{\arraystretch}{1.14}
\emergencystretch=3em
\setlist[itemize]{leftmargin=1.4em,itemsep=2pt,topsep=3pt}

\pagestyle{fancy}
\fancyhf{}
\fancyhead[L]{\small\color{muted}California FeatureGraph transfer}
\fancyhead[R]{\small\color{muted}23 July 2026}
\fancyfoot[C]{\small\color{muted}\thepage}
\renewcommand{\headrulewidth}{0.3pt}

\newcommand{\code}[1]{\texttt{#1}}
\newcommand{\model}[1]{\textbf{#1}}
\newcommand{\good}[1]{\textcolor{green}{\textbf{#1}}}
\newcommand{\warn}[1]{\textcolor{orange}{\textbf{#1}}}
\newcommand{\bad}[1]{\textcolor{red}{\textbf{#1}}}
\newcommand{\reportbox}[2]{%
  \begin{center}
  \fcolorbox{#1!45}{#1!6}{%
    \begin{minipage}{0.94\linewidth}
      \small #2
    \end{minipage}}
  \end{center}}

\title{\vspace{-1.0cm}
  \color{navy}\bfseries Cross-model transfer of evolved tabular FeatureGraphs\\[4pt]
  \large Eight graph sources, eight evaluators, five estimator seeds\\
  \normalsize California Housing}
\author{GigaEvo research report}
\date{23 July 2026}

\begin{document}
\maketitle
\vspace{-0.7cm}

\begin{abstract}
\noindent
We froze eight FeatureGraphs evolved against CatBoost, TabM, RealMLP,
TabICLv2, TabPFN~v3, LightGBM, XGBoost, and TabFM~1.0, then evaluated every
graph with every estimator on the same held-out California Housing split.
The complete matrix contains 64 graph--evaluator cells plus eight raw-feature
controls, each summarized over estimator seeds 0--4.  Fifty-seven engineered
cells beat the corresponding raw control on every matched seed.  The seven
exceptions reveal a sharp asymmetry: every graph selected by the older
evaluators hurts TabFM, while TabFM's compact seven-feature graph improves all
eight evaluators over raw.  CatBoost still wins five rows, TabPFN wins the two
non-CatBoost boosting rows, and TabFM's native graph wins the TabFM row.
Exact output-name overlap remains tiny (2.7\% median pairwise Jaccard).
All searches recover geography and household normalization, but TabFM alone
selects a compact, entirely unsupervised spatial graph rather than supervised
neighbourhood summaries.  Transfer is therefore broad but not universal:
the consuming model's feature prior matters.
\end{abstract}

\reportbox{blue}{
\textbf{\color{navy}Bottom line.}
Feature engineering learned against one model is often portable, but TabFM is
a clean counterexample to universality.  CatBoost remains the safest broad
donor for the seven older evaluators and TabPFN remains the strongest foreign
donor on average.  TabFM instead prefers its own narrow, unsupervised graph;
all seven foreign graphs degrade it.  This is a post-hoc audit on one fixed
test split, so these patterns are hypotheses to pre-register on new datasets,
not a license to select graphs on this test matrix.}

\section{Experimental question and protocol}

\subsection{What is transferred}

A FeatureGraph is executable feature code, not a fitted estimator.  Its JSON
topology and transforms are frozen, but aggregate nodes and the downstream
model are rebuilt inside each evaluation using only the permitted fitting
partition.  Thus matrix row \(i\), column \(j\) asks:

\begin{quote}
\emph{How well does evaluator \(i\) perform when given the exact feature
program selected by evaluator \(j\)'s evolution?}
\end{quote}

Rows are evaluators; columns are graph sources.  The raw column retains only
the eight California inputs.  The eight graph columns retain all raw inputs
and append the generated numerical columns.  No evolution, mutation model, or
memory model was called during the matrix run.

\subsection{Data and fitting protocol}

The frozen split contains 13,209 training rows, 3,303 validation rows, and
4,128 test rows.  The regression target is median house value in units of
\$100k.  The eight inputs are median income, house age, average rooms,
average bedrooms, population, average occupancy, latitude, and longitude.

For CatBoost, TabM, RealMLP, LightGBM, and XGBoost, validation selects the
training length under the model's fixed recipe; the graph, preprocessing, and
estimator are then rebuilt cleanly for the final train-plus-validation fit.
TabICL, TabPFN, and TabFM are frozen in-context estimators, so train plus
validation is their labeled context rather than an optimizer trajectory.  The
test labels are read only after prediction to compute RMSE and \(R^2\).

Each cell uses estimator seeds \(0,\ldots,4\).  Reported ``\(\pm\)'' values
are sample standard deviations (\(\mathrm{ddof}=1\)) across those five
estimator fits.  The graph and data split do not vary.  LightGBM and XGBoost
are deterministic under their fixed recipes, hence their zero seed SD.
CatBoost varies because its experiment-local constructor consumes the
estimator seed.

\subsection{Execution and completeness}

The final artifact contains 72 cells and 360 seed-level records.  The original
seven-model study contributed 56 cells.  The TabFM extension added its raw
control, eight evaluator-row cells, and seven donor-column cells---80 new
seed-level evaluations at revision \code{05e27851}.  Within the original
matrix, 42 previously missing cells were run at \code{26ab5fe3} and fourteen
protocol-compatible cells were reused.  GPU evaluators leased randomly from
four H100s through the shared cross-process allocator; tree evaluators used
CPU workers.  Every extension process exited successfully, every cell
contains exactly seeds 0--4, all metrics are finite, and graph hashes match
the frozen files.

\section{The transfer matrix}

\begin{figure}[H]
  \centering
  \includegraphics[width=\linewidth]{figures/rmse_matrix.pdf}
  \caption{\textbf{Absolute held-out RMSE.}  Each cell is mean \(\pm\) sample
  SD over five estimator seeds.  Dark solid outlines mark the minimum RMSE in
  each evaluator row; dashed orange outlines mark the graph evolved against
  that same evaluator.  Lower is better.}
  \label{fig:rmse}
\end{figure}

The strongest absolute result is TabFM consuming its native graph:
\(0.358948\pm0.000148\) RMSE and
\(0.901190\pm0.000081\) \(R^2\).  Absolute values across evaluator rows also
reflect model strength, so the row pattern is more informative:

\begin{itemize}
  \item \model{CatBoost's graph wins CatBoost, TabM, RealMLP, TabICL, and
  TabPFN.}  It is top-three for all seven older evaluators, but ranks fourth
  under TabFM.
  \item \model{TabPFN's graph wins LightGBM and XGBoost.}  It is second under
  CatBoost and TabPFN, and top-three in seven of eight rows.
  \item \model{TabFM's graph wins TabFM.}  Its \(0.358948\) RMSE is
  \(0.001019\) below raw, and \(0.006883\) below the best foreign graph.
  Native graphs otherwise win only CatBoost.  Champions were selected by
  evolutionary CV, whereas Figure~\ref{fig:rmse} is a post-freeze test audit
  and must not become a selection loop.
\end{itemize}

\begin{table}[H]
\centering
\scriptsize
\resizebox{\linewidth}{!}{%
\begin{tabular}{lrrrrrr}
\toprule
Evaluator & Raw RMSE & Native RMSE & Best graph & Best RMSE & Best \(R^2\) & Native rank \\
\midrule
\input{results/best_results_rows.tex}
\end{tabular}}
\caption{Raw, native, and row-best held-out results.  ``Best'' is descriptive
after observing the complete matrix; it is not a test-valid selection rule.}
\label{tab:best}
\end{table}

\begin{figure}[H]
  \centering
  \includegraphics[width=0.94\linewidth]{figures/improvement_matrix.pdf}
  \caption{\textbf{Relative RMSE reduction against each row's raw-feature
  control.}  The comparison is evaluator-local, so it removes the large
  difference in raw model strength.  Fifty-seven of 64 graph cells improve
  raw on all five matched seeds; the seven negative cells are every foreign
  graph consumed by TabFM.  Solid and dashed outlines have the same meaning
  as in Figure~\ref{fig:rmse}.}
  \label{fig:uplift}
\end{figure}

The largest graph effect is on RealMLP: CatBoost features cut its raw RMSE by
16.17\%, compared with 13.31\% for RealMLP's native graph.  Even the strongest
raw evaluator, TabPFN, benefits from every graph; its improvements are smaller
(0.67--4.17\%) because its raw RMSE is already \(0.378814\).  TabFM behaves
differently: its native graph improves raw by only 0.28\%, while foreign
graphs worsen RMSE by 1.63--4.58\%.  This reversal occurs on every matched
seed, not only in the five-seed mean.

\subsection{Foreign graph versus native graph}

The paired quantity below is
\(\mathrm{RMSE}_{\mathrm{native}}-\mathrm{RMSE}_{\mathrm{best\ foreign}}\);
positive values favor transfer.  Pairing uses the same estimator seed.

\begin{table}[H]
\centering
\small
\begin{tabular}{lrrr}
\toprule
Evaluator & Best foreign graph & Paired native--foreign RMSE & Foreign wins \\
\midrule
\input{results/native_foreign_rows.tex}
\end{tabular}
\caption{Native versus the best foreign source.  Seed counts are descriptive,
not a substitute for uncertainty across datasets or splits.}
\label{tab:paired}
\end{table}

The CatBoost donor advantage is large and seed-consistent for RealMLP
(\(+0.013170\pm0.002925\), five of five) and clear for TabM and TabICL.
For TabPFN the \(+0.000958\) mean advantage reverses on one seed, so it should
be treated as a near-tie.  CatBoost itself remains better with its own graph
than with the best foreign graph.  The deterministic LightGBM and XGBoost
gains from TabPFN are \(0.005911\) and \(0.011625\) RMSE respectively.
TabFM is the strongest native-specific case: its graph beats the best foreign
source (TabICL) by \(0.006883\pm0.000114\) on all five seeds.

\section{Who is a good feature donor?}

\begin{figure}[H]
  \centering
  \includegraphics[width=\linewidth]{figures/transfer_summary.pdf}
  \caption{\textbf{Donor and family-level transfer.}  Left: bars show each
  graph's mean RMSE reduction on the seven non-native evaluators; diamonds show
  its native reduction; bar colors denote boosting (blue), classical DL
  (orange), and foundation (purple).  Right: mean evaluator-local reduction
  for each source-family/evaluator-family block.  Family cells are descriptive
  averages over individual graph/evaluator pairs.}
  \label{fig:donors}
\end{figure}

\begin{table}[H]
\centering
\small
\begin{tabular}{llrrrr}
\toprule
Graph source & Family & Native gain & Foreign mean & Row wins & Top-three rows \\
\midrule
\input{results/donor_summary_rows.tex}
\end{tabular}
\caption{Donor summary, ordered by mean gain on the seven foreign evaluators.}
\label{tab:donors}
\end{table}

Two complementary donors emerge:

\begin{itemize}
  \item \model{CatBoost remains the broad donor for the original model set.}
  Its 8.80\% mean foreign gain is not the largest, but five row wins and seven
  top-three placements show low downside outside TabFM.
  \item \model{TabPFN is the high-upside compact donor.}  Its 9.47\% foreign
  mean is the largest.  It
  transfers exceptionally to all three boosting evaluators, winning LightGBM
  and XGBoost despite containing only 14 generated columns.
  \item \model{LightGBM is a strong consensus donor without a win.}  It is
  top-three in six rows and averages an 8.17\% foreign reduction.
  \item \model{TabFM is a weak but consistently positive foreign donor.}  Its
  graph improves all seven older evaluators on every seed, averaging 4.59\%,
  yet ranks last in each of those rows.  Its value is specific to TabFM,
  where it is the only engineered graph that beats raw.
\end{itemize}

At family level, boosting-origin graphs are strongest for classical DL
evaluators (12.08\% mean reduction), while foundation-origin graphs transfer
strongly to boosting (9.47\%).  Adding TabFM sharply lowers every mean for
foundation evaluators: boosting, classical-DL, and foundation graphs average
only 2.71\%, 1.49\%, and 2.47\% respectively, because the seven foreign
TabFM-row cells are negative.  Family averages therefore conceal the central
asymmetry and should be read alongside Figure~\ref{fig:uplift}.

\section{What each evolution run discovered}

\begin{figure}[H]
  \centering
  \includegraphics[width=\linewidth]{figures/graph_anatomy.pdf}
  \caption{\textbf{Semantic allocation and graph structure.}  The semantic
  taxonomy assigns each emitted column to exactly one interpretable family;
  it is a manual audit of graph code, not learned feature importance.  The
  structure panel distinguishes rowwise-node and aggregate-node outputs.
  Aggregate does not necessarily mean target-derived: it means the node has a
  fitted transform and may use fitting-set context.}
  \label{fig:anatomy}
\end{figure}

\begin{table}[H]
\centering
\small
\begin{tabular}{llrrrrr}
\toprule
Graph & Family & Nodes & Depth & Features & Rowwise outputs & Aggregate outputs \\
\midrule
\input{results/graph_structure_rows.tex}
\end{tabular}
\caption{Exact topology of the eight frozen graphs.  All eight raw columns are
retained in addition to the generated columns shown here.}
\label{tab:structure}
\end{table}

\subsection{CatBoost: broad, balanced, and compositional}

The five-node, depth-two graph emits 33 features.  Three independent rowwise
nodes create six coordinate rotations (30, 45, and 60 degrees), four
Haversine city distances plus two distance ratios, and five household or
demographic ratios.  An aggregate node adds target means at \(k=4,16,64\),
target dispersion, three density measures, and local income, age, room, and
occupancy summaries.  A dependent node then contrasts each row with those
neighborhood summaries and forms a micro-to-macro target ratio.

This is the most heterogeneous representation: geometry, household
normalization, local covariates, supervised context, and cross-scale
composition all receive capacity.  That balance is the most plausible
explanation for its five evaluator wins and top-three placement everywhere.

\subsection{TabM: a dense multi-scale neighborhood basis}

TabM also emits 33 features, but 28 come from one flat aggregate node.  At
\(k=5,15,50,150\) it repeats target mean and SD, local-income mean and SD,
spatial density, and income-to-neighborhood ratios and differences.  The five
rowwise features are LA/SF distance, bedroom share, income times house age,
and income per occupant.

The graph is semantically concentrated rather than small: four scales repeat
the same local-market hypothesis.  It performs well across all evaluators,
but its redundancy does not buy universal donor status; the broader CatBoost
and LightGBM representations beat it under TabM itself.

\subsection{RealMLP: smooth spatial bases plus model stacking}

RealMLP emits 17 features in three nodes.  Its large rowwise node combines
three household ratios, a 40.6-degree coordinate rotation, and distance plus a
50\,km radial-basis response for LA, SF, San Diego, San Jose, and Sacramento.
In parallel it creates an out-of-fold 15-neighbor spatial regressor.  A
dependent aggregate node fits an out-of-fold
\code{HistGradientBoostingRegressor} on raw and engineered inputs and emits
one stacked prediction.

No other graph uses a supervised model-stack output, and no other graph spends
ten columns on paired city-distance/RBF bases.  This makes the representation
distinctive and effective natively, but comparatively specialized in foreign
evaluators.

\subsection{TabICL: coastal gating and small modular nodes}

TabICL's six-node graph emits 15 columns.  It builds a coastal decay signal,
then uses it in coastal-income and inland-income log interactions.  Parallel
nodes add distances to SF, LA, San Diego, and Sacramento; three household
ratios; target KNN summaries at \(k=5,15,50\); and local income and occupancy
anomaly ratios.

Its unique contribution is not another city anchor but a soft
coastal-versus-inland gate.  Structurally it is the most modular graph, with
one explicit dependency and five largely independent mechanisms.

\subsection{TabPFN: compact robust neighborhood statistics}

TabPFN emits only 14 features from four independent nodes: LA/SF distances;
local income mean, income SD, and spatial density; three household ratios; and
both mean and median target summaries at \(k=4,16,64\).

The median summaries are unique among the eight graphs and give a robust
counterpart to local means.  This graph has neither coordinate rotations nor
a large local-covariate bank, yet it is the best foreign donor on average.
Its success argues that concise, multi-scale target context plus a few stable
rowwise ratios can be more portable than sheer width.

\subsection{LightGBM: a compact version of the CatBoost portfolio}

LightGBM emits 16 features in five nodes: two rotated coordinates, three
household ratios, LA/SF distances, and eight neighborhood features covering
local income, house age, density, target means, and target/income dispersion.
A dependent node divides household income by neighborhood income.

Its semantic allocation is closest to CatBoost's (cosine 0.91), but in a much
smaller graph.  This helps explain its six top-three placements: it retains
most of the balanced portfolio while avoiding many repeated scales.

\subsection{XGBoost: learned spatial bases and local covariate detail}

XGBoost emits 32 features through seven independent nodes.  Besides LA/SF
distance, a bedroom ratio, and a 45-degree coordinate rotation, it learns ten
KMeans centroid-distance features.  It adds target KNN means at
\(k=5,15,50\), twelve local income/house-age means and SDs over the same
scales, plus income per capita and rooms per person.

The ten centroid distances are unique: they form an unsupervised radial basis
over California instead of using only named cities.  The graph is rich, but
its large local-covariate and clustering allocation is less transferable than
the smaller TabPFN graph under XGBoost itself.

\subsection{TabFM: a narrow unsupervised spatial prior}

TabFM's four independent nodes emit only seven features: distance to San
Francisco and Los Angeles; bedrooms per room and occupants per room; latitude
and longitude rotated by 35 degrees; and mean distance to the ten nearest
training coordinates.  It is the only graph with no target-derived aggregate
or out-of-fold supervised predictor.  Its target transform standardizes on
the fitting partition and clips the inverse prediction to the documented
California target range.

The allocation is deliberately narrow in effect, though it was discovered
without such a prompt: six rowwise features plus one unsupervised spatial
density.  It gives TabFM a small but seed-consistent 0.28\% gain, whereas the
wider foreign graphs---especially those dominated by supervised
neighbourhood summaries---hurt TabFM.  Yet the same seven features improve
every older evaluator over raw, making this graph portable in one direction
and highly native-specific in the other.

\section{How similar are the feature sets?}

\begin{figure}[H]
  \centering
  \includegraphics[width=\linewidth]{figures/similarity_matrices.pdf}
  \caption{\textbf{Lexical difference versus semantic convergence.}  Left:
  Jaccard overlap of exact emitted column names.  Right: cosine similarity of
  the mutually exclusive semantic-count profiles in
  Figure~\ref{fig:anatomy}.  The diagonal is included as a scale check.
  Semantic similarity is an analyst-coded summary, not a performance metric.}
  \label{fig:similarity}
\end{figure}

Exact naming makes the graphs look almost unrelated: median off-diagonal
Jaccard is 2.7\%, mean is 3.7\%, and the maximum is 14.3\% for
RealMLP--TabICL.  Much of that disagreement is superficial:
\code{dist\_la}, \code{dist\_to\_la}, and a California-coordinate RBF can
encode the same underlying geography while sharing no string.

At the semantic level, two ideas appear in \emph{all eight} graphs:

\begin{enumerate}[leftmargin=1.6em,itemsep=2pt,topsep=3pt]
  \item explicit geographic proximity or a coordinate basis;
  \item household-scale normalization such as bedrooms per room or rooms per
  occupant.
\end{enumerate}

The seven older graphs additionally use leakage-probed spatial target context
at one or more scales (or an out-of-fold spatial predictor in RealMLP).
TabFM is the exception: it uses only unsupervised neighbour spacing.

The feature sets therefore differ more in \emph{allocation and implementation}
than in their central hypothesis.  CatBoost and LightGBM have cosine 0.91
because both balance coordinate transforms, household ratios, local
covariates, target summaries, and a dependent neighborhood contrast.
TabPFN and LightGBM also score 0.91 despite very different names.  TabFM is
semantically close to CatBoost (0.83) and RealMLP (0.78) because of their
geographic allocations, but close semantics do not prevent negative transfer
into TabFM.  RealMLP and TabM remain the least similar pair at 0.17: RealMLP
emphasizes smooth fixed anchors and two model predictions; TabM repeats
neighborhood statistics across four scales.

Semantic similarity alone does not determine transfer.  CatBoost's breadth
reduces mismatch among the original evaluators but does not help TabFM.
TabPFN's compact target mean/median portfolio transfers especially well to
trees, yet also hurts TabFM.  RealMLP's distinctive stacking is valuable to
RealMLP but not broadly dominant.  Graph size is likewise insufficient: the
14-feature TabPFN graph is the best foreign donor by mean, while TabFM's
seven-feature graph is last in every foreign row but uniquely successful in
its native row.

\section{Interpretation and next experiments}

\reportbox{green}{
\textbf{\color{navy}Working interpretation.}
The searches did not discover eight unrelated recipes.  They discovered a
shared geographic representation of California, then allocated finite graph
capacity according to the evaluator's inductive bias.  Transfer succeeds
among the original evaluators because much of that representation is useful
across models.  TabFM exposes the boundary: its frozen in-context prior
prefers a compact unsupervised augmentation and is degraded by every wider
foreign representation.}

The matrix suggests three pre-registered follow-ups:

\begin{itemize}
  \item \textbf{Conditional donor test.}  Freeze the CatBoost graph before
  touching test labels on new regression datasets, then test separately on
  trained estimators and frozen in-context estimators.  The present matrix
  rejects an unconditional universal-donor claim.
  \item \textbf{Compact donor test.}  Repeat with the TabPFN graph.  Its 14
  columns give a clean width-controlled challenger and test whether
  mean-plus-median neighborhood summaries are broadly tree-friendly.
  \item \textbf{TabFM width and supervision ablation.}  Starting from raw,
  add TabFM's seven features, then matched-width subsets of foreign graphs,
  then target-derived neighbourhood summaries.  Freeze the ablation under CV
  before test to separate width, supervision, and semantic mismatch.
  \item \textbf{Semantic union and pruning.}  Build unions by semantic family,
  not by concatenating every column.  Add one mechanism at a time under CV
  (for example CatBoost's coordinate/composition blocks plus TabPFN's median
  summaries), remove correlated duplicates, freeze once, then evaluate on a
  new test split or dataset.  A raw eight-graph union would confound feature
  quality with input width and estimator parameter count.
\end{itemize}

The report does \emph{not} recommend selecting the lowest cell in this matrix
for a production model.  The matrix has already consumed the test split.  Its
proper use is to formulate smaller hypotheses and carry them to untouched
datasets.

\section{Limitations and provenance}

\begin{itemize}
  \item \textbf{One dataset and one split.}  Seed SD measures estimator
  randomness only; it is not uncertainty over samples, splits, tasks, or
  evolution trajectories.
  \item \textbf{Post-hoc matrix.}  Original champions were frozen without
  using test labels, but cross-matrix row winners are observed on test and
  cannot be used as an unbiased selection rule.
  \item \textbf{Search campaigns are not a causal estimator ablation.}  The
  six newest graphs came from matched 100-mutation Memory~V2 campaigns;
  CatBoost and TabM came from earlier completed campaigns.  Prompts, model
  explanations, memory histories, and source revisions are not perfectly
  identical.  The matrix supports claims about consuming frozen graphs, not
  about which estimator causes evolution to invent a mechanism.
  \item \textbf{Graph width and model capacity vary.}  Generated width ranges
  from 14 to 33.  For neural models, width can change parameterization as well
  as representation.
  \item \textbf{Historical compatible cells are reused.}  The complete
  per-cell revision is stored in \code{combined\_matrix.json}.  Four
  CatBoost/TabM cells use \code{4d4a6cfc}; two RealMLP raw/native cells use
  \code{680cbed}; eight other raw/native cells use \code{9d23221}; the 42 new
  original cells use \code{26ab5fe3}; and the 16 TabFM-extension cells use
  \code{05e27851}.  The relevant frozen evaluator paths were
  compatibility-audited, but revision boundaries remain visible rather than
  being hidden.
  \item \textbf{Foundation checkpoint licensing.}  TabPFN~v3 uses the
  Prior-Labs checkpoint and TabFM uses the Google TabFM checkpoint under
  their respective non-commercial terms.
\end{itemize}

The tested environment used Python~3.12, PyTorch~2.13, TabM~0.0.3,
PyTabKit~1.7.3, TabICL~2.1.1, TabPFN~8.1.0, TabFM~1.0.1,
CatBoost~1.2.10, LightGBM~4.6.0, and XGBoost~2.1.4.

\subsection{Frozen graph identities}

\begin{table}[H]
\centering
\small
\begin{tabular}{ll}
\toprule
Graph & SHA-256 \\
\midrule
Raw      & \code{a08f6e89573958de\ldots b6b699c3} \\
CatBoost & \code{ea1a75f553598033\ldots 1a8ccfb7} \\
TabM     & \code{7d859c823c0af111\ldots d3cb565e} \\
RealMLP  & \code{4d239c46c4471857\ldots effc8bb5} \\
TabICL   & \code{ab8df0bdf2563868\ldots d6c9dd4c} \\
TabPFN   & \code{8755fc2919c9663d\ldots ba75336} \\
LightGBM & \code{a275bde3203176db\ldots 9e5716fb} \\
XGBoost  & \code{9847d9e5e31e61da\ldots 062c331} \\
TabFM    & \code{95e44f3454fa04a2\ldots e5058b271} \\
\bottomrule
\end{tabular}
\caption{Byte hashes of the self-contained graph files shipped with this
report.}
\end{table}

\section*{Artifact map}

The report directory is self-contained:

\begin{itemize}
  \item \code{REPORT.pdf} and \code{REPORT.tex}: rendered and source report;
  \item \code{graphs/}: the raw control and eight exact executable graphs;
  \item \code{results/combined\_matrix.json}: all 360 seed-level records,
  summaries, revisions, graph hashes, and paired raw deltas;
  \item \code{results/*.csv}: RMSE, uplift, rank, donor, family, paired, and
  similarity tables; and
  \item \code{make\_figures.py} and \code{figures/}: reproducible plotting
  source plus vector PDF and high-resolution PNG assets.
\end{itemize}

\vfill
\begin{center}
\small\color{muted}
Generated from frozen local artifacts; no LLM or evolution credits were used
for the transfer matrix.
\end{center}

\end{document}
